% Compile with: latexmk -pdf main
\documentclass[11pt]{article}

\usepackage[margin=1in]{geometry}
\usepackage{amsmath,amssymb,amsthm}
\usepackage[T1]{fontenc}
\usepackage[utf8]{inputenc}
\usepackage{lmodern}
\usepackage{newunicodechar}
\usepackage{fancyvrb}
\usepackage{xcolor}
\usepackage[hidelinks]{hyperref}

% Unicode characters appearing in Lean code
\newunicodechar{‖}{\ensuremath{\|}}
\newunicodechar{ℝ}{\ensuremath{\mathbb{R}}}
\newunicodechar{ℚ}{\ensuremath{\mathbb{Q}}}
\newunicodechar{→}{\ensuremath{\to}}
\newunicodechar{∀}{\ensuremath{\forall}}
\newunicodechar{∃}{\ensuremath{\exists}}
\newunicodechar{∈}{\ensuremath{\in}}
\newunicodechar{∧}{\ensuremath{\wedge}}
\newunicodechar{∩}{\ensuremath{\cap}}
\newunicodechar{≤}{\ensuremath{\leq}}
\newunicodechar{β}{\ensuremath{\beta}}
\newunicodechar{ε}{\ensuremath{\varepsilon}}
\newunicodechar{θ}{\ensuremath{\theta}}
\newunicodechar{π}{\ensuremath{\pi}}
\newunicodechar{ρ}{\ensuremath{\rho}}
\newunicodechar{Ω}{\ensuremath{\Omega}}
\newunicodechar{ᶠ}{\ensuremath{{}^{f}}}
\newunicodechar{𝓝}{\ensuremath{\mathcal{N}}}
\newunicodechar{₀}{\ensuremath{{}_0}}
\newunicodechar{√}{\ensuremath{\surd}}
\newunicodechar{×}{\ensuremath{\times}}
\newunicodechar{∫}{\ensuremath{\int}}

\DefineVerbatimEnvironment{lean}{Verbatim}{fontsize=\small,frame=leftline,framerule=0.6pt,
  rulecolor=\color{gray!60},xleftmargin=1em}

\newtheorem{theorem}{Theorem}
\newtheorem{lemma}[theorem]{Lemma}
\theoremstyle{remark}
\newtheorem{remark}[theorem]{Remark}

\title{Closed Near-Circular Orbits Force a Power Law:\\
A Formal Proof of the First Half of Bertrand's Theorem in Lean~4}
\author{Nichlas Löfdahl \and Claude (Anthropic)\thanks{AI system (model Claude Opus 5.5), developed by Anthropic.
Claude designed the formalization, wrote the Lean proofs (in part via parallel sub-agents) and drafted this
manuscript in collaboration with the first author.}}
\date{October 2026}

\begin{document}
\sloppy
\maketitle

\begin{abstract}
Bertrand's theorem (1873) states that the only central forces for which every bounded orbit is
closed are the inverse-square law and Hooke's law. All standard proofs share a first step. If
every orbit near each stable circular orbit is closed, then the near-circular apsidal angle
$\pi/\sqrt{3+rV''/V'}$ must be a constant rational multiple of $\pi$, so the force is a power law
$V'(r)=C\,r^{\beta^2-3}$ with $\beta\in\mathbb{Q}_{>0}$. We give a machine-checked proof of this
step in Lean~4 with Mathlib, with the apsidal angle defined as an honest singular integral over
the classically allowed region. The analytic core is a one-dimensional Morse lemma, proved via
L'Hôpital's rule. A trigonometric substitution in Morse coordinates then
turns the apsidal integral into a regular integral. Together they show that the apsidal angle is
continuous in the energy up to the circular orbit, where it equals $\pi/\sqrt{U''}$. To our
knowledge this is the first formalization of any part of Bertrand's theorem.
\end{abstract}

\section{Introduction}

Among all central forces, why do only gravity and the ideal spring produce closed orbits? Bertrand's
answer~\cite{bertrand1873} is a classic of mechanics, reproduced in Goldstein~\cite{goldstein} and
Arnold~\cite{arnold}. Rigorous treatments are surprisingly delicate. Arnold's sketch hides
nontrivial analysis in its limits. Féjoz and Kaczmarek~\cite{fejoz2004} give a careful local proof
via Birkhoff invariants, and many elementary proofs have since appeared (e.g.~\cite{chin2015}).

Most proofs proceed in two stages:
\begin{enumerate}
  \item[(I)] \emph{Near-circular orbits.} If all orbits close to stable circular orbits are
  closed, the force is a power law $F\propto r^{\beta^2-3}$ with rational $\beta$.
  \item[(II)] \emph{Power laws.} Among power laws, only $\beta=1$ (Kepler) and $\beta=2$ (harmonic
  oscillator) have all bounded orbits closed.
\end{enumerate}
This paper formalizes stage (I) completely in Lean~4~\cite{mathlib2020}. Stage (II) and the link
between closed trajectories and apsidal angles are left for future work
(\S\ref{sec:discussion}). The development is part of the repository
\url{https://github.com/Crotalus/ai-formalized} (folder \texttt{bertrand}). It builds with Lean
\texttt{v4.35.0} and Mathlib commit \texttt{56a0058}.

\section{Setting}

A particle of unit mass moves in a central potential $V(r)$ with angular momentum $M>0$. In the
variable $x=M/r$, Binet's equation states that the radial motion, parametrized by the polar angle
$\theta$, is a one-dimensional oscillation $\tfrac12 x'(\theta)^2+U(x)=E$ in the effective potential
\[
  U(x)=\frac{x^2}{2}+V\!\left(\frac{M}{x}\right).
\]
The \emph{apsidal angle}, the angle from pericentre to apocentre, is the half-period of this
oscillation:
\[
  \Phi(E)=\int_{\{U<E\}_{x_0}}\frac{dx}{\sqrt{2(E-U(x))}},
\]
where $\{U<E\}_{x_0}$ is the connected component, containing the bottom $x_0$ of the well, of the
classically allowed region. An orbit closes after finitely many revolutions exactly when
$\Phi/\pi\in\mathbb{Q}$. A circular orbit of radius $r$ has $M^2=r^3V'(r)$ and sits at the
critical point $x_0=M/r$ of $U$, with
\begin{equation}\label{eq:Upp}
  U''(x_0)=3+\frac{rV''(r)}{V'(r)} ,
\end{equation}
and it is stable when this is positive. In Lean:

\begin{lean}
def apsidalAngle (U : ℝ → ℝ) (Ω : Set ℝ) (x₀ E : ℝ) : ℝ :=
  ∫ x in connectedComponentIn (Ω ∩ {x | U x < E}) x₀, 1 / √(2 * (E - U x))

def effPot (V : ℝ → ℝ) (M x : ℝ) : ℝ := x ^ 2 / 2 + V (M / x)
def angMom (V : ℝ → ℝ) (r : ℝ) : ℝ := √(r ^ 3 * deriv V r)
\end{lean}

The integrand is unbounded at the turning points, so $\Phi$ is a genuinely singular (Lebesgue)
integral. We do not regularize it in the definition.

\section{Main theorem}

\begin{theorem}[Bertrand, part I]\label{thm:main}
Let $0<a<b$ and let $V$ be $C^3$ on $(a,b)$ with $V'>0$ and $3V'+rV''>0$ there (every circular
orbit of radius in $(a,b)$ exists and is stable). Suppose that for every $r\in(a,b)$, every orbit
with the angular momentum of the circular orbit of radius $r$ and energy slightly above it is closed,
that is, $\Phi(E)\in\pi\mathbb{Q}$. Then there is a rational $\beta>0$ with $3+rV''/V'=\beta^2$ on
$(a,b)$, and $V'(r)=C\,r^{\beta^2-3}$ for some $C>0$.
\end{theorem}

\begin{lean}
theorem forceLaw_of_closed_nearCircular {V : ℝ → ℝ} {a b : ℝ} (ha : 0 < a) (hab : a < b)
    (hV : ContDiffOn ℝ 3 V (Ioo a b))
    (hpos : ∀ r ∈ Ioo a b, 0 < deriv V r)
    (hstab : ∀ r ∈ Ioo a b, 0 < 3 * deriv V r + r * iteratedDeriv 2 V r)
    (hclosed : ∀ r ∈ Ioo a b, ∀ᶠ E in 𝓝[>] (circularEnergy V r),
      ∃ q : ℚ, orbitApsidalAngle V a b r E = q * π) :
    ∃ β : ℚ, 0 < β ∧ (∃ C > 0, ∀ r ∈ Ioo a b, deriv V r = C * r ^ ((β : ℝ) ^ 2 - 3)) ∧
      ∀ r ∈ Ioo a b, 3 + r * iteratedDeriv 2 V r / deriv V r = (β : ℝ) ^ 2
\end{lean}

The orbit's allowed region is taken inside the radii $(a,b)$, i.e.\ $\Omega=(M/b,M/a)$ in the
$x$ variable.

\paragraph{Proof outline.} Fix $r$ and the well $U$ of its circular orbit.
\begin{enumerate}
  \item \emph{Morse chart} (\S\ref{sec:morse}). There is a $C^1$ map $X$ with $X(0)=x_0$, $X'>0$
  and $U(X(s))=U(x_0)+s^2/2$ near $0$, with $X'(0)=1/\sqrt{U''(x_0)}$.
  \item \emph{Regularization} (\S\ref{sec:chart}). For $0<\varepsilon$ small,
  $\Phi(U(x_0)+\varepsilon)=\int_{-\pi/2}^{\pi/2}X'(\sqrt{2\varepsilon}\sin\theta)\,d\theta=:F(\varepsilon)$.
  The function $F$ is continuous on $[0,\rho^2/2)$, with $F(0)=\pi X'(0)=\pi/\sqrt{U''(x_0)}$.
  \item \emph{Rationality.} By hypothesis $F(\varepsilon)/\pi\in\mathbb{Q}$ for small
  $\varepsilon>0$. A continuous rational-valued function on an interval is constant, since
  otherwise its image would contain an irrational number. So $F$ is constant near $0^+$, and by
  continuity $F(0)\in\pi\mathbb{Q}$. Hence $\sqrt{U''(x_0)}\in\mathbb{Q}$.
  \item \emph{Power law.} By \eqref{eq:Upp}, $r\mapsto\sqrt{3+rV''/V'}$ is continuous on $(a,b)$
  and rational-valued, hence a constant $\beta$. Then $rV''=(\beta^2-3)V'$, so
  $(V'(r)\,r^{3-\beta^2})'=0$ and $V'=C\,r^{\beta^2-3}$.
\end{enumerate}

Step 3 is where continuity of $\Phi$ \emph{in the energy} is indispensable. A function with values
in $\pi\mathbb{Q}$ can perfectly well have an irrational limit. This is why the near-circular
formula $\Phi\to\pi/\sqrt{U''}$, which textbooks derive by linearization, does not by itself
prove the step.

\section{The one-dimensional Morse lemma}\label{sec:morse}

\begin{lean}
theorem exists_morseChart {U : ℝ → ℝ} {Ω : Set ℝ} {x₀ : ℝ} (hΩ : IsOpen Ω) (hx₀ : x₀ ∈ Ω)
    (hU : ContDiffOn ℝ 3 U Ω) (h1 : deriv U x₀ = 0) (h2 : 0 < iteratedDeriv 2 U x₀) :
    ∃ c : MorseChart U Ω x₀, deriv c.X 0 = 1 / √(iteratedDeriv 2 U x₀)
\end{lean}

The usual proof writes $U(x)-U(x_0)=(x-x_0)^2g(x)$ with
$g(x)=\int_0^1(1-t)U''(x_0+t(x-x_0))\,dt$, and needs differentiation under the integral sign. The
formal proof avoids this. It sets $q(x)=(U(x)-U(x_0))/(x-x_0)^2$ and
$s(x)=(x-x_0)\sqrt{2q(x)}$, and uses Mathlib's L'Hôpital rule (\texttt{lhopital\_zero\_nhdsNE}) to
get $q(x)\to U''(x_0)/2$. Hence $s'(x_0)=\sqrt{U''(x_0)}$. Away from $x_0$, differentiating
$s^2=2(U-U(x_0))$ gives $s'=U'/s$, which tends to $\sqrt{U''(x_0)}$ as well. So $s$ is $C^1$ with
positive derivative near $x_0$, and the local inverse given by Mathlib's inverse function theorem
(\texttt{HasStrictFDerivAt.toOpenPartialHomeomorph}) is the chart $X$.

\section{The apsidal angle in a Morse chart}\label{sec:chart}

\begin{lean}
theorem MorseChart.apsidalAngle_eq (c : MorseChart U Ω x₀) {ε : ℝ} (hε : 0 < ε)
    (hερ : ε < c.ρ ^ 2 / 2) :
    apsidalAngle U Ω x₀ (U x₀ + ε) =
      ∫ θ in -(π / 2)..π / 2, deriv c.X (√(2 * ε) * sin θ)
\end{lean}

Let $\sigma=\sqrt{2\varepsilon}$. First, the allowed component equals $X((-\sigma,\sigma))$. This
image is connected, lies in $\{U<E\}$ and contains $x_0$. Conversely, the component is an interval
around $x_0$ that cannot contain the turning points $X(\pm\sigma)$, where $U=E$. Second, the map
$\theta\mapsto X(\sigma\sin\theta)$ is a $C^1$ bijection from $(-\pi/2,\pi/2)$ onto that
component. On it $2(E-U)=\sigma^2\cos^2\theta$, so the Jacobian $X'(\sigma\sin\theta)\sigma\cos\theta$
cancels the singularity exactly. The change of variables uses Mathlib's Lebesgue version
\texttt{integral\_image\_eq\_integral\_abs\_deriv\_smul}, which requires neither continuity nor
integrability of the singular integrand. Continuity of
$F(\varepsilon)=\int X'(\sqrt{2\varepsilon}\sin\theta)\,d\theta$ on $[0,\rho^2/2)$ follows from
continuity of parametric integrals with jointly continuous integrand.

\section{Discussion}\label{sec:discussion}

\paragraph{Size and checking.} The development consists of five files, about 750 lines in all. It
has no \texttt{sorry}, and \texttt{\#print axioms} on the main theorem reports only
\texttt{propext}, \texttt{Classical.choice} and \texttt{Quot.sound}. The proof was developed in
parallel. We first fixed a small shared interface: the definition of \texttt{apsidalAngle} and a
\texttt{MorseChart} structure. Separate agents then proved the Morse lemma and the chart
computation against it, while the physics layer and the assembly were written in parallel.

\paragraph{Stage II.} With $V'=Cr^{\beta^2-3}$ on all of $(0,\infty)$ and all bounded orbits
closed, the same continuity-plus-rationality argument, applied over the connected range of bound
energies, shows that $\Phi(E)=\pi/\beta$ for every bound energy. Comparing with the limit
$E\to\infty$ for $\beta^2>2$ (where $\Phi\to\pi/2$) and $E\to0^-$ for $\beta^2<2$ (where
$\Phi\to\pi/\beta^2$) leaves $\beta=2$ and $\beta=1$. Formalizing this needs a global Morse chart
for a single-well potential and two asymptotic evaluations of the apsidal integral. The
regularized integral of \S\ref{sec:chart} is designed to make both tractable.

\paragraph{Trajectories.} We take as given that an orbit closes exactly when
$\Phi\in\pi\mathbb{Q}$. Deriving this from Newton's equations (Binet's equation and periodicity of
the radial motion) would connect the result to Physlib's central-force library.

\paragraph{Author contributions.} N.L.\ chose the problem, directed the work and reviewed the
results. Claude, an AI system developed by Anthropic, designed the formal statements and proof
architecture, wrote the Lean proofs (delegating the Morse lemma and the chart computation to
sub-agents) and drafted the manuscript. All statements were checked by the Lean kernel. As an AI
system cannot take responsibility for a publication, N.L.\ is accountable for its content.

\bibliographystyle{plain}
\bibliography{refs}

\end{document}
